\paragraph{Selecting the acquisition protocol.}
The interface, budget unit and operational preference must be specified before
an acquisition guarantee is used. A reference implementation selects among
the supplied fixed, adaptive and full-range certificates for the 29 cases.
It checks permitted
queries, scalar-reading and information-round caps, transport grouping, and
the requested preference; a missing guarantee triggers escalation rather than
a relaxed target. This is an integration of the preceding results, not a new
optimization theorem. A frozen offline evaluation retains all 29 external
cases, the exact endpoint total, the $0.01E$ stored-readout allowance and the
unchanged half-width target.

\begin{table}[t]
\centering\small
\caption{Protocol selection under four explicitly declared interfaces. The
restricted interfaces are benchmark assumptions; the public API permits bulk
retrieval. All rows use the same 29 cases and uncertainty model.}
\label{tab:protocol-selection}
\begin{tabular}{p{0.40\linewidth}cp{0.36\linewidth}}
\hline
Interface and preference & Certified & Selected protocol \\
\hline
Range available; fewest requests & 29/29 & Bulk range in all 29 \\
Three readings, two rounds; fewest readings & 29/29 & Adaptive two in 22; fixed three in seven \\
Two readings, two rounds & 22/29 & Adaptive two in 22; escalate seven \\
Three readings, one round & 29/29 & Fixed three in all 29 \\
\hline
\end{tabular}
\end{table}

Unrestricted range retrieval is preferred on the actual public interface.
Its five internal values contain each case's certified three-query subset,
so information monotonicity preserves the precision guarantee under bounded
stored-readout errors; it does not imply zero uncertainty about true increments.
The restricted three-reading rule reserves 65 scalar observations in 51
information rounds, whereas the one-round-per-case rule reserves 87 in 29
rounds. These are distinct resources. HTTP counts are left unspecified for
the hypothetical restricted interfaces. In particular, an interface offering
a single native multipoint request can favor fixed-three under request
priority while favoring adaptive-two under reading priority.

The preceding three-task API timings support bulk selection on that service;
they are not measurements on this 29-case sensitivity subset. The supplement exports all
116 case--interface decisions, exclusions and certificate hashes, with an
offline callback dispatcher and separate verification. The implementation returns an auditable acquisition plan: a supported
case, its measurement assumptions, selected query protocol, hard resource
allowances and the certificate proving its precision. If no supplied plan
meets the interface, it returns escalation. This catalogue decision alone
does not prove that every possible policy fails; the seven deterministic
two-query exclusions additionally have the exact lower witnesses described
above.

The operational wrapper binds each supplied certificate to its saved model
pair, exact coefficients, meter, API-grid block, retained endpoints and source
hashes; readout envelopes with a different task, timestamp, feed or unit are
rejected. For example, the system~1 ridge--MLP comparison on 6 July 2025,
12:00--18:00 UTC, uses boundary~2 followed by boundary~1 or~4 under the
restricted two-reading rule, while its public-interface plan selects bulk
retrieval. These are checks against accidental task mismatch and a worked
planning example; measured API trials and constructed dispatch tests remain
separate, and the wrapper does not authenticate a provider or calibrate a meter.
